% chapters/02/section-03.tex
\section{Ruang, Besaran, dan Representasi Keadaan}
\label{sec:ruang-keadaan}

Setelah waktu, unsur kedua yang harus disepakati adalah ruang: di mana objek
berada, seberapa besar ukurannya, dan bagaimana keadaan permainan dinyatakan
di dalam program.

\subsection{Sistem Koordinat Layar dan Kartesius}

Layar komputer menyusun gambar dari titik-titik yang tersusun dalam baris dan
kolom. Titik kiri-atas dipakai sebagai titik asal, dan sumbu menurun ke bawah
seiring bertambahnya nomor baris. Cara ini berbeda dari sistem kartesius yang
Anda pelajari di sekolah, tempat titik asal biasanya berada di kiri-bawah dan
sumbu menanjak ke atas. Kedua sistem diperlihatkan pada
Gambar~\ref{fig:koordinat}.

\begin{figure}[htbp]
\centering
% chapters/02/figures/koordinat.tikz
% Dua sistem koordinat: koordinat layar (asal kiri-atas, sumbu y ke bawah)
% dan koordinat kartesius (asal kiri-bawah, sumbu y ke atas).
\begin{tikzpicture}[font=\scriptsize]
  % --- Koordinat layar: asal di pojok kiri-atas, y menurun ---
  \draw[ubPrimary, thick, fill=ubAlt] (0,0) rectangle (4.6,3.2);
  \draw[-{Stealth[length=2mm]}, ubAccent, thick] (0,3.2) -- (1.4,3.2)
    node[anchor=south west] {$x$};
  \draw[-{Stealth[length=2mm]}, ubAccent, thick] (0,3.2) -- (0,1.8)
    node[anchor=east] {$y$};
  \fill[ubRed] (0,3.2) circle (1.6pt);
  \node[anchor=north west, ubRed, inner sep=3pt] at (0.14,3.08) {asal $(0,0)$};
  \node[anchor=north, ubPrimary, font=\footnotesize\bfseries] at (2.3,-0.25)
    {koordinat layar};
  \node[anchor=north, font=\scriptsize\itshape, text width=48mm] at (2.3,-0.7)
    {titik asal di kiri-atas; $y$ bertambah ke bawah};

  % --- Koordinat kartesius: asal di kiri-bawah, y menaik ---
  \begin{scope}[xshift=7.6cm]
    \draw[ubPrimary, thick, fill=ubAlt] (0,0) rectangle (4.6,3.2);
    \draw[-{Stealth[length=2mm]}, ubAccent, thick] (1.2,0.7) -- (2.6,0.7)
      node[anchor=west] {$x$};
    \draw[-{Stealth[length=2mm]}, ubAccent, thick] (1.2,0.7) -- (1.2,2.1)
      node[anchor=south] {$y$};
    \fill[ubRed] (1.2,0.7) circle (1.6pt);
    \node[anchor=north west, ubRed, inner sep=3pt] at (1.34,0.62)
      {asal $(0,0)$};
    \node[anchor=north, ubPrimary, font=\footnotesize\bfseries] at (2.3,-0.25)
      {koordinat kartesius};
    \node[anchor=north, font=\scriptsize\itshape, text width=48mm] at (2.3,-0.7)
      {titik asal di kiri-bawah; $y$ bertambah ke atas};
  \end{scope}
\end{tikzpicture}

\caption{Sistem koordinat layar dan sistem koordinat kartesius}
\label{fig:koordinat}
\end{figure}

Perbedaan ini bukan persoalan sepele. Nilai positif pada sumbu tegak berarti
``ke bawah'' ketika Anda bekerja dengan koordinat layar, sehingga gravitasi
yang seharusnya menarik objek ke bawah justru dinyatakan sebagai penambahan
nilai, bukan pengurangan. Kebingungan tentang arah sering menjadi penyebab
pertama mengapa objek pada program latihan bergerak ke arah yang berlawanan
dengan dugaan penulisnya.

Mesin permainan tertentu memakai sistem ketiga, yaitu koordinat dunia tiga
dimensi dengan sumbu tegak yang dapat dipilih penataannya. Yang penting bukan
sistem mana yang benar, melainkan kesadaran bahwa sistem yang dipakai harus
diketahui secara pasti pada setiap tahap pekerjaan.

\subsection{Skala dan Satuan}

Persoalan berikutnya adalah ukuran. Berapa piksel panjang satu meter? Berapa
besar sebuah karakter dibandingkan lebar layar? Pertanyaan semacam ini
diistilahkan sebagai \istilah{metrics} dan dijawab lebih dahulu pada tahap
perancangan level, karena angka yang dipilih menentukan rasa permainan
\cite{unity-leveldesign}.

Menetapkan skala secara sengaja membawa tiga manfaat. Pertama, ukuran objek
menjadi konsisten, sehingga karakter setinggi dua petak tidak tiba-tiba
berubah menjadi setinggi dua puluh piksel pada layar lain. Kedua, kecepatan
gerak dapat dibandingkan secara bermakna; kecepatan berjalan dan berlari dapat
dinyatakan sebagai perbandingan angka yang sama pada setiap level. Ketiga,
perhitungan fisika menjadi lebih mudah diatur karena rentang nilainya masuk
akal, tidak terlalu kecil maupun terlalu besar.

\begin{catatabox}
Pada Pygame, satuan yang paling wajar dipakai adalah piksel, sebab koordinat
pada layar memang dinyatakan dalam piksel. Pada Unity, satuan dasar yang
dipakai fisika adalah meter, dan satu satuan koordinat sama dengan satu meter.
Perbedaan ini akan Anda rasakan pada Bab~XII ketika berpindah dari satu
perkakas ke perkakas lainnya.
\end{catatabox}

\subsection{Besaran Turunan: Posisi, Kecepatan, dan Percepatan}

Keadaan gerak sebuah objek dinyatakan oleh tiga besaran yang saling
berkaitan. Besaran-besaran ini bukan sekadar nama, melainkan tingkatan
perubahan terhadap waktu.

\begin{description}
  \item[Posisi] Letak objek pada suatu saat, dinyatakan sebagai pasangan atau
        rangkap tiga angka.
  \item[Kecepatan] Laju perubahan posisi terhadap waktu. Nilainya menyatakan
        seberapa cepat, sedangkan arahnya menyatakan ke mana.
  \item[Percepatan] Laju perubahan kecepatan terhadap waktu. Gravitasi adalah
        contoh percepatan yang tetap besarnya dan mengarah ke bawah.
\end{description}

Hubungan antarketiganya bersifat berulang: kecepatan mengubah posisi, dan
percepatan mengubah kecepatan. Urutan pembaruan inilah yang menentukan hasil
simulasi, sebagaimana diuraikan pada Bagian~\ref{sec:simulasi-diskret}. Angka
yang menyatakan besaran tersebut dikelompokkan menjadi vektor, yaitu kumpulan
nilai yang diperlakukan sebagai satu kesatuan karena menyatakan besaran
berarah. Pada Bab~V Anda akan memakai jenis data bawaan Pygame untuk keperluan
itu \cite{pygame-docs}.

\subsection{Ruang Batas dan Tabrakan}

Untuk menentukan apakah dua objek bersentuhan, program tidak memeriksa setiap
titik pada gambar objek. Cara yang dipakai lebih sederhana: setiap objek
dilingkupi oleh persegi pembatas yang sisi-sisinya sejajar dengan sumbu
koordinat, disebut \istilah{axis-aligned bounding box} (AABB). Dua objek
dianggap bersentuhan apabila persegi pembatasnya beririsan.

Pendekatan ini dipilih karena murah dihitung dan cukup untuk sebagian besar
kebutuhan. Harga yang dibayar adalah ketelitian: dua objek yang bentuknya
tidak mengisi seluruh persegi pembatasnya dapat terdeteksi bersentuhan
walaupun gambar keduanya tidak benar-benar bersinggungan. Apabila ketelitian
itu diperlukan, tersedia cara yang lebih cermat dengan memeriksa piksel yang
benar-benar terisi \cite{pygame-sprite-mask}. Pertimbangan memilih di antara
keduanya dibahas pada Bab~VII, sedangkan contoh kajian penerapannya pada
permainan dapat dilihat pada \mbox{[R23]}.

\subsection{Representasi Keadaan Permainan}

Keadaan permainan adalah kumpulan seluruh nilai yang menentukan apa yang
sedang terjadi. Nilai-nilai itu terbagi menjadi dua golongan.

\begin{itemize}
  \item \textbf{Keadaan yang terus berubah} dengan nilai pecahan: posisi,
        kecepatan, sisa waktu tunggu, dan kemajuan animasi.
  \item \textbf{Keadaan yang berubah secara terputus} dengan nilai tertentu
        saja: keadaan NPC, apakah pemain sedang melompat, fase permainan
        (bermain, jeda, selesai), skor, dan jumlah nyawa.
\end{itemize}

Membedakan keduanya bermanfaat. Golongan pertama menuntut perhitungan yang
cermat dan penyajian yang halus, sedangkan golongan kedua menentukan aturan
permainan dan lebih mudah diperiksa. Sebagian besar kesalahan logika
permainan berasal dari golongan kedua: keadaan yang seharusnya sudah berubah
ternyata belum, atau berubah dua kali padahal seharusnya sekali.

Satu anjuran praktis yang berlaku sepanjang buku ini: nyatakan keadaan
permainan secara terpusat dan sejelas mungkin. Apabila Anda menemukan diri
Anda menyimpan satu informasi yang sama di dua tempat berbeda, hampir dapat
dipastikan kedua salinan itu akan saling bertentangan pada suatu saat. Pada
Bab~IX hal ini akan Anda jumpai kembali ketika menyusun mesin keadaan, tempat
aturan ``satu keadaan aktif pada satu waktu'' mencegah percabangan logika yang
rumit.
